%% IOS Press FAIA — iosart2x.cls
%% JURIX 2025: International Conference on Legal Knowledge and Information Systems
%%
%% Hybrid Sparse-Dense Retrieval for EU Climate Law:
%%   Evaluating BM25, BGE-M3, and Reciprocal Rank Fusion
%%   on a Statute-Level Benchmark
%%
\documentclass{iosart2x}

\usepackage[T1]{fontenc}
\usepackage{natbib}           % \citet / \citep author-year citations
\usepackage{amsmath,amssymb}
\usepackage{booktabs}
\usepackage{graphicx}
\usepackage{microtype}
\usepackage[ruled,vlined,linesnumbered]{algorithm2e}
\usepackage[hidelinks]{hyperref}
\usepackage{xspace}
\usepackage{multirow}

%% Convenience macros
\newcommand{\HR}[1]{\text{HR}@{#1}}
\newcommand{\MRR}[1]{\text{MRR}@{#1}}
\newcommand{\PH}{\text{[PLACEHOLDER]}\xspace}
\newcommand{\PHCI}[1]{\text{[PH]~(CI: [PH, PH])}}
\newcommand{\bgem}{BGE-M3\xspace}
\newcommand{\rrf}{RRF\xspace}

% ──────────────────────────────────────────────────────────────────────────────
\begin{document}

\begin{frontmatter}

\title{Hybrid Sparse-Dense Retrieval for EU Climate Law:\\
       Evaluating BM25, \bgem{}and Reciprocal Rank Fusion\\
       on a Statute-Level Benchmark}

\author[A]{\inits{R.S.}\fnm{Rameez} \snm{Shafat}}

\address[A]{%
  [Institution]\\
  \email{[email address]}%
}

\begin{abstract}
Grounded responses in legal AI applications depend critically on
first-stage retrieval: a system that fails to surface the correct
statutory provision cannot produce a verifiable, citation-supported
answer.  We evaluate three retrieval configurations---BM25 lexical
search, \bgem dense retrieval (BAAI/bge-m3, 1024\,dim), and a hybrid
combining both via Reciprocal Rank Fusion (\rrf, $k{=}60$)---on a
manually constructed benchmark of 71 queries over EU Climate Law.
The corpus comprises 1,156 article-level provisions from 72 EU
instruments extracted from the EU CELLAR database via SPARQL and
Formex~XML.  We measure retrieval effectiveness using $\HR{1}$,
$\HR{5}$, $\HR{10}$, and $\MRR{k}$ with Wilson 95\,\% confidence
intervals, and complement these with reliability (per-query hit
patterns) and efficiency (index size and query latency) analyses.
To the best of our knowledge, this is the first systematic comparative
study of hybrid sparse-dense retrieval on EU statutory text, and the
first publicly available statute-level retrieval benchmark for EU
Climate Law.  Experimental results are reported as \PH{}pending
completion of the evaluation runs.  The findings directly inform the
design of first-stage retrieval components in grounded legal AI
systems where source attribution to specific CELEX instruments and
articles is mandatory.
\end{abstract}

\begin{keyword}
EU Climate Law\sep legal information retrieval\sep hybrid retrieval\sep
BM25\sep BGE-M3\sep reciprocal rank fusion\sep retrieval-augmented
generation\sep statute-level benchmark
\end{keyword}

\end{frontmatter}

% ──────────────────────────────────────────────────────────────────────────────
\section{Introduction}
\label{sec:intro}

Consider the following failure: a legal analyst queries a BM25-only system
with \textit{``What monitoring obligations apply to aircraft operators under
the EU Emissions Trading System?''}\  The system returns top-ranked
provisions from Regulation~(EU)~2015/757 on maritime emissions
monitoring---\textsc{celex} \texttt{32015R0757}---because the terms
\textit{monitoring}, \textit{operators}, and \textit{obligations} achieve a
high BM25 score against that instrument's preamble and operative articles.
The correct provisions, Articles~14--15 of Directive~2003/87/EC
(\texttt{32003L0087}), rank below position~10 owing to the
aviation-specific vocabulary (\textit{tonne-kilometre data},
\textit{verified emissions report}) absent from the query.  The downstream
language model, receiving a context that contains no mention of aviation
or Article~14, cannot produce a grounded answer and either hallucinates a
response or correctly invokes an insufficient-context guard.

This failure illustrates why first-stage retrieval quality is the
rate-limiting factor for grounded legal AI applications.
Retrieval-augmented generation~(RAG) pipelines for statute interpretation
require that the retriever surface not merely topically similar text, but
the legally binding provisions that \emph{answer} the query.  In EU
Climate Law, this requirement is compounded by legislative density
(72~enacted instruments), cross-instrument definitional dependencies
(e.g.,\ the CBAM Regulation cross-references ETS allowance concepts
defined in Directive~2003/87/EC), highly technical and
domain-specific vocabulary (\textit{free allocation},
\textit{carbon leakage}, \textit{corrigendum}), and the urgent policy
context that makes error-free source attribution a professional and
legal necessity.

Sparse lexical retrieval (BM25) captures exact statutory terminology
but is susceptible to vocabulary mismatch~\cite{robertson2009bm25}.
Dense semantic retrieval (\bgem) generalises beyond exact matches but
may conflate legally distinct concepts sharing surface
similarity~\cite{chen2024bgem3}.  Hybrid fusion via Reciprocal Rank
Fusion (\rrf) combines both lists without requiring score normalisation
or labelled training data~\cite{cormack2009rrf}, making it particularly
suitable for a domain where annotated retrieval pairs are scarce.

\paragraph{Research Questions.}

\begin{description}
  \item[\textbf{RQ1.}] Does hybrid sparse-dense retrieval
    (BM25\,+\,\bgem{}with \rrf $k{=}60$) improve the effectiveness and
    reliability of first-stage grounding for EU Climate Law compared
    with BM25-only and \bgem-only retrieval?
  \item[\textbf{RQ2.}] Is first-stage hybrid retrieval sufficient to
    support accurate and domain-relevant grounded responses in EU
    Climate Law, given retrieval coverage, hard-query failure rates,
    and efficiency overhead?
\end{description}

\paragraph{Contributions.}  This paper makes four contributions:

\begin{enumerate}
  \item A 71-query manually verified benchmark for EU Climate Law
    statute-level retrieval with binary CELEX-level relevance
    judgements (\S\ref{sec:eval-bench}).
  \item A reproducible article-level corpus of 1,156 provisions from
    72 EU climate instruments, extracted from the EU CELLAR SPARQL
    endpoint via a documented two-phase pipeline
    (\S\ref{sec:corpus}).
  \item A three-way comparative evaluation of BM25, \bgem{}, and
    BM25+\bgem{} \rrf across effectiveness, reliability, and
    efficiency dimensions (\S\ref{sec:results}).
  \item An analysis of which first-stage retrieval failure modes most
    threaten reliable source attribution in domain-specific LLM
    deployments, with concrete design implications
    (\S\ref{sec:error}).
\end{enumerate}

% ──────────────────────────────────────────────────────────────────────────────
\section{Related Work}
\label{sec:related}

\subsection{Lexical Retrieval in Legal IR}
\label{sec:related-lexical}

BM25 remains the dominant first-stage lexical retrieval model owing to
its probabilistic grounding and robust performance without training
data~\cite{robertson2009bm25}.  In legal retrieval, exact match over
statutory vocabulary is particularly valuable: defined terms such as
\textit{installation}, \textit{allowance}, and \textit{tonne-kilometre}
carry precise legal meanings that synonymy-based generalisation may
distort.  \citet{mori2023} examine statute-level retrieval for legal
question answering and find that lexical signals remain strong for
definitional and obligation-type queries where the relevant provision
shares key vocabulary with the question.  The NOWJ team at
COLIEE~\cite{nowjcoliee2023} demonstrate that well-tuned BM25 baselines
are competitive on statutory entailment tasks, establishing a high bar
for learned alternatives.  The fundamental limitation of pure lexical
retrieval---term mismatch between user query vocabulary and statutory
vocabulary---motivates the dense component evaluated in this paper.

\subsection{Dense Retrieval and Legal Embeddings}
\label{sec:related-dense}

Dense retrieval encodes both queries and documents into a shared
embedding space, enabling semantic matching beyond exact term
overlap~\cite{chen2024bgem3}.  The \bgem model (\texttt{BAAI/bge-m3})
achieves leading performance on the MTEB multilingual retrieval
benchmark through self-knowledge distillation and multi-granularity
encoding~\cite{chen2024bgem3}; critically for this work, its 8,192-token
context window accommodates long legislative articles without arbitrary
truncation for the majority of the corpus.  \citet{rayo2024} evaluate
dense retrieval for regulatory compliance queries and identify
distributional shift between general-domain pre-training corpora and
EU statutory text as a factor limiting performance.  \citet{legalrag2024}
integrate dense retrieval into a legal QA pipeline and report that
semantic matching reduces hard failures on queries whose vocabulary
diverges from the document collection.  A shared limitation is that
dense models may retrieve thematically related but legally distinct
instruments---a failure mode we characterise in \S\ref{sec:error}.

\subsection{Sparse Learned Representations}
\label{sec:related-splade}

SPLADE~\cite{formal2021splade} learns sparse representations by applying
a log-saturation activation over BERT vocabulary logits, producing
inverted-index-compatible vectors that encode both lexical and expanded
signals.  The subsequent SPLADE-v2 extension~\cite{formal2022splade2}
improves effectiveness through distillation from a dense teacher and
hard-negative sampling, substantially narrowing the gap between sparse
learned models and dense retrievers on standard BEIR benchmarks.
However, both variants require domain-specific fine-tuning to be
effective: without in-domain annotated pairs, SPLADE models exhibit
vocabulary hallucination---expanding queries towards general-domain
concepts absent from the target corpus.  No SPLADE model trained on
EU statutory text currently exists.  This gap justifies using
terminology-preserving BM25 rather than a sparse learned model as
the lexical component, and flags SPLADE fine-tuning on EU law as
a productive direction for future work.

\subsection{Hybrid Retrieval and Rank Fusion}
\label{sec:related-hybrid}

Hybrid retrieval combining sparse and dense ranked lists via Reciprocal
Rank Fusion was shown by~\citet{cormack2009rrf} to outperform individual
rank learning methods and Condorcet fusion on TREC collections without
requiring score normalisation.  \citet{hyparag2024} apply hybrid
parameter-adaptive retrieval to AI legal and policy applications,
demonstrating that fusion over sparse and dense lists reduces failure
rates on complex regulatory queries compared with either component
alone.  \citet{lrage2024} evaluate multiple hybrid retrieval
configurations in a legal RAG benchmark framework and find that rank
fusion consistently improves reliability over single-retriever baselines,
though the magnitude varies with query type and domain specificity.
A recurring finding across these systems is that the independence
assumption underlying \rrf---that ranking errors of each retriever
are partially uncorrelated---holds in legal IR because lexical and
semantic failure modes are empirically distinct, a claim we evaluate
in \S\ref{sec:reliability}.

\subsection{Legal RAG Systems}
\label{sec:related-legalrag}

Several systems address the end-to-end legal RAG pipeline.
\citet{cbrrag2023} augment retrieval with case-based reasoning,
using structural similarity between retrieved precedents and the
current query as a re-ranking signal.  \citet{lexrag2024} propose
LexRAG, combining lexical retrieval with generation for legal
reasoning tasks, and report that lexical precision is a stronger
predictor of answer correctness than semantic similarity on statutory
interpretation queries.  \citet{lexdrafter2024} apply RAG to
legislative drafting assistance, finding that article-level chunking
aligns with drafting practice better than fixed-token windows.
\citet{legalstructurerag2024} exploit legal document structure
(recitals, chapters, articles, sub-paragraphs) as a retrieval signal,
reporting gains on tasks where the correct provision is embedded
in a specific structural position.  A consistent gap across this
body of work is the absence of systematic evaluation on EU statutory
text; most systems focus on common-law jurisdictions (US, UK)
or case law retrieval, with EU civil-law legislation---which has
distinct structural and terminological properties---largely
unexplored.

\subsection{Retrieval Reliability and Source Attribution}
\label{sec:related-reliability}

\citet{reliableretrievallegal2024} identify retrieval reliability---the
consistency with which a system surfaces the correct provision across
paraphrased variants of the same query---as a distinct evaluation
dimension from single-query effectiveness, and argue that unreliable
retrieval is a principal driver of hallucination in grounded legal AI.
\citet{reasoninglegal2024} frame legal retrieval as a reasoning
problem in which the first-stage retriever must resolve implicit
definitional dependencies between instruments before a chain-of-thought
generator can be applied.  Both works underscore that for legally
consequential applications---where every claim in an answer must be
attributed to a specific, verifiable statutory provision---single-point
IR metrics understate the practical risk introduced by retrievers with
uneven per-query coverage.  We address this gap directly through the
reliability dimension in \S\ref{sec:eval-reliability}.

\subsection{Evaluation Frameworks for Legal Retrieval}
\label{sec:related-eval}

The COLIEE benchmark~\cite{nowjcoliee2023} is the most widely used
legal retrieval evaluation, covering statute retrieval and legal
entailment tasks primarily over the Japanese Civil Code and Canadian
case law.  \citet{lrage2024} propose a framework for evaluating legal
RAG systems that decomposes performance into retrieval coverage,
attribution accuracy, and answer faithfulness.  HyPA-RAG~\cite{hyparag2024}
includes an evaluation component for legal and policy question answering.
None of these frameworks provides a benchmark for EU Climate Law
statute retrieval, where the target collection spans 72 instruments in
the CELLAR repository and queries must resolve cross-instrument
definitional dependencies.  The benchmark constructed in this paper
fills this gap and is designed to be directly reproducible from the
publicly available CELLAR SPARQL endpoint.

% ──────────────────────────────────────────────────────────────────────────────
\section{Research Gap and Contributions}
\label{sec:gap}

The Related Work establishes three specific gaps that this paper
addresses:

\begin{description}
  \item[Corpus gap.]  No published collection of EU Climate Law
    provisions at article-level granularity, extracted from CELLAR and
    validated against a Pydantic schema, is currently available.  The
    corpus assembled here (1,156 provisions, 72 instruments) provides
    the first reproducible baseline for this domain.

  \item[Benchmark gap.]  No manually constructed query benchmark for
    EU Climate Law statute-level retrieval exists.  Existing benchmarks
    (COLIEE, LRAGE) address different jurisdictions and document types.
    The 71-query benchmark developed here uses binary CELEX-level
    relevance judgements calibrated to EU legislative practice.

  \item[Evaluation gap.]  Although hybrid retrieval has been studied in
    general and legal IR, no prior work has systematically compared
    BM25, \bgem, and \rrf specifically on EU statutory text under a
    shared evaluation protocol that reports Wilson confidence
    intervals, per-query reliability, and efficiency simultaneously.
\end{description}

These gaps jointly define a contribution that cannot be addressed by
re-running any existing system on any existing benchmark.

% ──────────────────────────────────────────────────────────────────────────────
\section{Methodology}
\label{sec:method}

Figure~\ref{fig:architecture} gives an overview of the full pipeline.
The quantitative evaluation covers the shaded retrieval components;
the generation stage is described for completeness but evaluated
only qualitatively.

\begin{figure}[t]
  \centering
  \fbox{\rule{0pt}{4.5cm}\rule{0.88\linewidth}{0pt}}
  \caption{System architecture.  A query is processed concurrently by
    BM25 and \bgem; the two ranked lists of 20 candidates each are
    merged by \rrf to produce a fused top-5 context.  The fused
    context is passed to Llama\,3.3-70B (Ollama, local inference) for
    strictly-grounded answer generation, with every claim required to
    cite \texttt{[CELEX\_ID\,---\,Article\,N]}.  Quantitative
    evaluation covers only the retrieval components (shaded region).}
  \label{fig:architecture}
\end{figure}

\subsection{Corpus Construction}
\label{sec:corpus}

The corpus is extracted from the EU CELLAR repository using a
documented two-phase SPARQL pipeline.

\paragraph{Phase 1---Batch SPARQL.}  A SPARQL query against the
CELLAR endpoint retrieves all work URIs associated with 21 EuroVoc
climate and energy concepts, yielding 194 unique CELEX identifiers
for Directives and Regulations.

\paragraph{Phase 2---Per-document manifestation retrieval.}
For each work URI, a second SPARQL query resolves the available
Formex~XML (\texttt{fmx4}) manifestation URL; where no Formex
manifestation exists, an HTML fallback is used.  Documents are
fetched with retry and exponential back-off and cached locally.

\paragraph{Parsing.}  Formex~XML is parsed with \texttt{lxml}.
Articles are identified by the \texttt{<ARTICLE>} element; article
text is reconstructed from \texttt{<ALINEA>} leaf nodes to avoid
duplication of numbered sub-paragraphs inherited from parent
elements.  Each article is serialised to a JSON record
\texttt{\{celex\_id, doc\_type, article\_number, article\_text,
cross\_references, concept\_ids\}} and validated against a Pydantic
schema before inclusion; records containing placeholder text are
discarded.

\paragraph{Corpus statistics.}  The final corpus contains
\textbf{1,156} article-level provisions from \textbf{72} CELEX
instruments (Directives and Regulations), covering ETS, CBAM,
Taxonomy, LULUCF, F-Gas, Maritime~MRV, ESR, the European Climate Law,
Energy Efficiency, CSDDD, and the Governance Regulation.  Mean article
length is 2,103~characters; the distribution is heavily right-skewed
(minimum: 49~characters; maximum: 137,713~characters).  The maximum
article exceeds \bgem's 8,192-token context window; truncation
behaviour for this extreme tail is acknowledged as a limitation in
\S\ref{sec:validity}.

\paragraph{Why article-level chunking.}  EU legislative articles are
drafted as semantically self-contained units: each article establishes
one obligation, one definition, or one procedural rule.
Article-level chunking avoids the arbitrary window-boundary decisions
of token-based chunking and produces retrieval units that align with
how lawyers reason about statute---specifically, the citation format
\texttt{[CELEX\_ID\,---\,Article\,N]} that downstream grounded
generation requires.

\subsection{Sparse Retrieval---BM25}
\label{sec:sparse}

The sparse retriever implements BM25Okapi (rank-bm25\,0.2.2) with
$k_1{=}1.5$ and $b{=}0.75$.  A custom legal tokeniser applies the
regex pattern \texttt{[a-zA-Z0-9](?:[a-zA-Z0-9\textbackslash-]*[a-zA-Z0-9])?}
before lower-casing.  This preserves hyphenated statutory terms
(\textit{net-zero}, \textit{CBAM-certificate}, \textit{cross-border})
as single tokens.  No stemming and no stopword removal are applied:
EU legal definitions are definitionally precise in ways that stemming
distorts---\textit{emission} (singular) and \textit{emissions} (plural)
are used with distinct intent in ETS legislation, and removing
stopwords can break fixed legal phrases such as \textit{free of charge}
and \textit{not later than}.

The index is built by scoring all 1,156 documents with BM25Okapi and
persisted via pickle.  At query time, BM25 scores all documents and
returns the top~20 ranked results as \texttt{RetrievedResult} objects
with article, score, and rank fields.

Lexical precision is a direct grounding asset: queries phrased in
statutory language will find exact-match provisions that are the
direct citation target.

\subsection{Dense Retrieval---BGE-M3}
\label{sec:dense}

The dense retriever uses \texttt{BAAI/bge-m3} loaded via the
\texttt{sentence-transformers} library~\cite{reimers2019sbert,chen2024bgem3}.
The model produces 1,024-dimensional vectors.

\paragraph{Model selection rationale.}  \bgem was selected for three
reasons: (1)~it is fully open-weight and requires no API key, making
results reproducible without cost; (2)~its 8,192-token context window
accommodates the vast majority of EU articles, which average 2,103
characters (${\approx}525$~tokens); (3)~it achieves state-of-the-art
performance on the MTEB multilingual retrieval benchmark, providing a
strong semantic baseline without domain-specific fine-tuning.

\paragraph{Indexing.}  Documents are encoded with
\texttt{encode\_corpus(texts, normalize\_embeddings=True,
batch\_size=64)}.  Corpus vectors are L2-normalised and loaded into a
FAISS \texttt{IndexFlatIP}; inner product on normalised vectors is
equivalent to cosine similarity.

\paragraph{Retrieval.}  The query is encoded with
\texttt{encode\_queries([query], normalize\_embeddings=True)},
producing a single normalised vector.  FAISS returns the top~20
nearest neighbours by inner product as exact nearest-neighbour search.

\paragraph{Grounding implication.}  Semantic matching bridges
vocabulary gaps between user query phrasing and statutory terminology,
recovering relevant provisions that BM25 misses due to term mismatch.

\subsection{Reciprocal Rank Fusion}
\label{sec:rrf}

Retrieval lists from BM25 and \bgem are merged via Reciprocal Rank
Fusion.  For each document $d$, the fused score is:

\begin{equation}
  \label{eq:rrf}
  \mathrm{RRF}(d) = \sum_{r \,\in\, \{\mathcal{R}_s,\, \mathcal{R}_d\}}
    \frac{1}{k + \mathrm{rank}_r(d)}
\end{equation}

\noindent where $k{=}60$ is the smoothing constant~\cite{cormack2009rrf}
and $\mathrm{rank}_r(d)$ is the 1-based rank of document $d$ in
retriever list $r$ (or $\infty$ if $d$ is absent from $r$).  Documents
are deduplicated by the composite key
\texttt{celex\_id::article\_number} before fusion; each unique article
is scored once.  The top~5 fused results ($k_f{=}5$) are returned as
\texttt{FusedResult} objects carrying the RRF score, the individual
dense and sparse ranks, and the article text.

\paragraph{Why RRF over learned fusion.}  Two properties make \rrf
appropriate here.  First, BM25 scores are unbounded positive reals
while cosine similarities lie in $[-1,\,1]$; combining raw scores
requires normalisation assumptions that may be violated.  \rrf
operates on ranks only and is therefore scale-invariant~\cite{cormack2009rrf}.
Second, no labelled relevance pairs for EU Climate Law retrieval exist
in any public training set; learned fusion weights require supervised
signal that is unavailable.

\paragraph{Concurrent execution.}  Both retrievers are invoked
concurrently via \texttt{ThreadPoolExecutor(max\_workers=2)}.  If one
retriever raises an exception, the controller logs a warning and
continues with the remaining list; if both fail, the exception
propagates.

\paragraph{Grounding hypothesis.}  Because BM25 and \bgem fail on
largely disjoint query subsets (vocabulary mismatch vs.\ semantic
drift), fusion is hypothesised to provide more consistent first-stage
coverage, reducing the fraction of queries for which no relevant
provision appears in the top-5 context window.  RQ1 tests this
hypothesis empirically.

\subsection{Algorithm}
\label{sec:algo}

Algorithm~\ref{alg:hybrid} formalises the full pipeline.

\begin{algorithm}[t]
\DontPrintSemicolon
\caption{Hybrid Sparse-Dense Retrieval Pipeline}
\label{alg:hybrid}
\KwIn{Query $q$;\; per-retriever limit $k_r{=}20$;\; fusion output
  $k_f{=}5$;\; RRF constant $k{=}60$}
\KwOut{Fused list $\mathcal{R}_f$,\; $|\mathcal{R}_f|{=}k_f$}
\BlankLine
\tcp{Step 1 — concurrent retrieval (ThreadPoolExecutor, max\_workers=2)}
$\mathcal{R}_s,\;\mathcal{R}_d \;\leftarrow\;
  \textbf{parallel}\;\{\,
    \textsc{BM25.Retrieve}(q,\,k_r),\;\;
    \textsc{BGE\text{-}M3.Retrieve}(q,\,k_r)
  \,\}$\;
\BlankLine
\tcp{Step 2 — RRF fusion}
$\mathcal{D} \leftarrow \textsc{Deduplicate}(\mathcal{R}_s \cup \mathcal{R}_d)$
  \tcp*{key: \texttt{celex\_id::article\_number}}
\ForEach{$d \in \mathcal{D}$}{
  $\mathrm{rrf}(d) \;\leftarrow\;
    \displaystyle\sum_{r\,\in\,\{\mathcal{R}_s,\,\mathcal{R}_d\}}
    \frac{1}{k + \mathrm{rank}_r(d)}$\;
}
\BlankLine
\tcp{Step 3 — return top-$k_f$ by RRF score}
\Return $\;\textsc{TopK}\!\left(\mathcal{D},\;k_f;\;\text{key}=\mathrm{rrf}\right)$\;
\end{algorithm}

\subsection{Evaluation Benchmark Construction}
\label{sec:eval-bench}

The benchmark comprises 71 manually constructed queries covering the
eleven EU Climate Law instrument families in the corpus.  Each query
is phrased as a natural-language legal question of the type encountered
in regulatory compliance, policy analysis, or statutory interpretation:
for example, \textit{``What monitoring obligations apply to aircraft
operators under the EU ETS?''} or \textit{``What are the free
allocation benchmarks for industrial installations under ETS Phase~4?''}

\paragraph{Relevance judgements.}  Relevance is binary at the CELEX
instrument level (one or two relevant documents per query).  A CELEX
ID is judged relevant if the instrument contains the provision or
provisions that directly and materially answer the query; thematically
related instruments that do not contain the operative provision are
judged not relevant.

\paragraph{Corrigendum normalisation.}  CELEX identifiers carrying a
corrigendum suffix (e.g., \texttt{32003L0087R(02)}) are matched
against gold identifiers after stripping the \texttt{R(...)} suffix,
so that corrigendum-form identifiers in the corpus correctly match
base-form identifiers in the benchmark.

\paragraph{Statistical note.}  At $n{=}71$ queries the Wilson 95\,\%
confidence interval width for $\HR{5}$ is approximately $\pm 0.06$--$0.12$
depending on the observed rate (Equation~\ref{eq:wilson}), which is
sufficient to distinguish systems separated by $\geq 0.15$ in
$\HR{5}$.

\paragraph{Why document-level retrieval.}  Legal citations refer to
instruments and articles, not sub-article passages.  Document-level
(CELEX-level) matching is therefore the appropriate proxy for grounded
answerability: if the correct instrument is not retrieved, the correct
article cannot appear in the LLM's context.

% ──────────────────────────────────────────────────────────────────────────────
\section{Experimental Setup}
\label{sec:setup}

\paragraph{Hardware.}  [Specify machine used: CPU, RAM, GPU if applicable.]

\paragraph{Software.}
\begin{itemize}
  \item Python~[version]
  \item \texttt{sentence-transformers}~$\geq$3.3.0
  \item \texttt{faiss-cpu}~1.14.2
  \item \texttt{rank-bm25}~0.2.2
  \item \texttt{numpy}~2.4.6
  \item \texttt{pydantic}~2.13.4
  \item Llama~3.3-70B-Instruct via Ollama~[version] (qualitative demo only)
\end{itemize}

All library versions are pinned in \texttt{requirements.txt}.  The
\bgem model weight hash is recorded at download time to ensure
reproducibility.

\paragraph{Index build.}
The BM25 index is built by fitting \texttt{BM25Okapi} over all 1,156
tokenised article texts and serialised via \texttt{pickle}.
The dense FAISS index is built by embedding all articles in batches
of~64 with \texttt{encode\_corpus}, L2-normalising the resulting
$1156 \times 1024$ matrix, and writing to disk via
\texttt{faiss.write\_index}.  Build times are reported in
Table~\ref{tab:efficiency}.

\paragraph{Query latency measurement.}
For each retrieval condition, we execute all 71 gold queries in five
independent runs and report median and 95th-percentile latency.
Times include tokenisation, score computation, and result assembly
but exclude index loading from disk.

% ──────────────────────────────────────────────────────────────────────────────
\section{Evaluation Framework}
\label{sec:eval}

\subsection{Dimension 1: Effectiveness}
\label{sec:eval-effectiveness}

\paragraph{Hit Rate at $k$.}  $\HR{k}$ is the fraction of queries for
which at least one relevant CELEX identifier appears in the top-$k$
results:

\begin{equation}
  \label{eq:hitrate}
  \HR{k} = \frac{1}{|Q|}\sum_{q \in Q}
    \mathbf{1}\!\left[\mathrm{rank}_q \leq k\right]
\end{equation}

\noindent where $\mathrm{rank}_q$ is the rank of the highest-ranked
relevant document for query $q$, and $|Q|{=}71$.  We report
$k \in \{1, 5, 10\}$.

\paragraph{Mean Reciprocal Rank.}  $\MRR{k}$ is the mean of the
reciprocal rank of the first relevant document:

\begin{equation}
  \label{eq:mrr}
  \MRR{k} = \frac{1}{|Q|}\sum_{q \in Q}
    \frac{1}{\mathrm{rank}_q}\cdot
    \mathbf{1}\!\left[\mathrm{rank}_q \leq k\right]
\end{equation}

We report $k \in \{5, 10\}$.

\paragraph{Note on nDCG.}  The benchmark provides binary relevance
judgements.  nDCG with binary labels yields a DCG that differs from
MRR only in how sub-optimal ranks are weighted; it adds no analytical
information beyond what $\HR{k}$ and $\MRR{k}$ already capture, and
its use would misrepresent the evaluation's signal.  We therefore omit
nDCG.

\paragraph{Confidence intervals.}  Wilson 95\,\% confidence intervals
are reported for all $\HR{k}$ values:

\begin{equation}
  \label{eq:wilson}
  \mathrm{CI}_{95\%}(\hat{p}) =
    \frac{\hat{p} + \frac{z^2}{2n}
    \;\pm\;
    z\sqrt{\frac{\hat{p}(1-\hat{p})}{n} + \frac{z^2}{4n^2}}}
    {1 + \frac{z^2}{n}}
\end{equation}

\noindent where $\hat{p}$ is the observed Hit Rate, $n{=}71$, and
$z{=}1.96$.

\paragraph{Statistical significance.}  For pairwise comparison of
$\HR{5}$ between two systems, we apply McNemar's test on the
per-query hit/miss binary vectors:
$\chi^2 = (b-c)^2/(b+c)$, where $b$ and $c$ count queries on which
only the first or only the second system hits, respectively.  Results
are reported in Table~\ref{tab:effectiveness}.

\subsection{Dimension 2: Reliability}
\label{sec:eval-reliability}

Per-query hit patterns across the three conditions are tabulated in
Table~\ref{tab:reliability}.  We report: (i)~queries where all three
systems hit; (ii)~queries where the hybrid strictly improves over
\emph{both} baselines; (iii)~queries where the hybrid regresses versus
at least one baseline; (iv)~queries that all three systems fail
(\emph{hard queries}).  Hard queries define the ceiling for first-stage
retrieval performance under the current corpus and are analysed in
\S\ref{sec:error}.

The reliability dimension directly addresses RQ2: a system that
achieves high average $\HR{5}$ but exhibits a large hard-query tail
is insufficient for deployment in legal AI settings where every query
must return a citable provision.

\subsection{Dimension 3: Efficiency}
\label{sec:eval-efficiency}

Table~\ref{tab:efficiency} reports index size on disk (MB), index
build time (seconds), median query latency (ms), and 95th-percentile
query latency (ms) for each condition.  We also report the overhead
ratio of hybrid latency to the faster single-retriever latency,
quantifying the cost of concurrent fusion.

% ──────────────────────────────────────────────────────────────────────────────
\section{Results}
\label{sec:results}

Table~\ref{tab:effectiveness} presents retrieval effectiveness.
Table~\ref{tab:efficiency} reports efficiency.
Table~\ref{tab:reliability} breaks down per-query hit patterns.
Table~\ref{tab:ablation} reports an ablation on the \rrf $k$ constant.

\begin{table}[t]
\caption{Retrieval Effectiveness ($n{=}71$ queries).
  Wilson 95\,\% confidence intervals in parentheses for $\HR{k}$.
  $\dagger$~McNemar's test vs.\ BM25-only; $\ddagger$~vs.\ \bgem-only
  ($p$-value reported).}
\label{tab:effectiveness}
\centering
\small
\begin{tabular}{@{}lllllll@{}}
\toprule
System
  & $\HR{1}$
  & $\HR{5}$
  & $\HR{10}$
  & $\MRR{5}$
  & $\MRR{10}$
  & McNemar $p$ \\
\midrule
BM25-only
  & \PH
  & \PH{} (\PH, \PH)
  & \PH{} (\PH, \PH)
  & \PH
  & \PH
  & --- \\[2pt]
\bgem-only
  & \PH
  & \PH{} (\PH, \PH)
  & \PH{} (\PH, \PH)
  & \PH
  & \PH
  & --- \\[2pt]
BM25+\bgem{} RRF
  & \PH
  & \PH{} (\PH, \PH)
  & \PH{} (\PH, \PH)
  & \PH
  & \PH
  & $\dagger$\PH, $\ddagger$\PH \\
\bottomrule
\end{tabular}
\end{table}

\begin{table}[t]
\caption{Efficiency.  Build time and latency measured on [hardware].
  Latency: median and P95 over 5 runs $\times$ 71 queries.
  Overhead ratio: hybrid median latency $\div$ faster single-retriever
  median latency.}
\label{tab:efficiency}
\centering
\small
\begin{tabular}{@{}lllll@{}}
\toprule
System & Index size & Build time & Median lat. & P95 lat. \\
       & (MB)       & (s)        & (ms)        & (ms)     \\
\midrule
BM25-only          & \PH & \PH & \PH & \PH \\
\bgem-only         & \PH & \PH & \PH & \PH \\
BM25+\bgem{} RRF   & \PH & \PH & \PH & \PH \\
\midrule
Overhead ratio     & --- & --- & \PH & \PH \\
\bottomrule
\end{tabular}
\end{table}

\begin{table}[t]
\caption{Reliability breakdown ($n{=}71$ queries).
  ``Hybrid strictly improves'' = hybrid hits; both baselines miss.
  ``Hybrid regresses'' = hybrid misses; at least one baseline hits.
  ``All miss'' = hard queries.}
\label{tab:reliability}
\centering
\small
\begin{tabular}{@{}lll@{}}
\toprule
Category & Queries ($n$) & Fraction \\
\midrule
All three systems hit                      & \PH & \PH \\
Hybrid strictly improves over both bases   & \PH & \PH \\
Hybrid regresses vs.\ $\geq$1 baseline    & \PH & \PH \\
All three systems miss (hard queries)      & \PH & \PH \\
\bottomrule
\end{tabular}
\end{table}

\begin{table}[t]
\caption{Ablation: sensitivity of hybrid \rrf to the $k$ constant.
  $\HR{5}$ and $\MRR{5}$ on the full 71-query benchmark.
  Baseline: $k{=}60$.}
\label{tab:ablation}
\centering
\small
\begin{tabular}{@{}lll@{}}
\toprule
$k$   & $\HR{5}$ & $\MRR{5}$ \\
\midrule
10  & \PH & \PH \\
30  & \PH & \PH \\
60  & \PH & \PH \\
100 & \PH & \PH \\
\bottomrule
\end{tabular}
\end{table}

% ──────────────────────────────────────────────────────────────────────────────
\section{Discussion}
\label{sec:discussion}

\paragraph{RQ1: Hybrid effectiveness and reliability.}
Table~\ref{tab:effectiveness} shows that BM25+\bgem{} \rrf achieves
$\HR{5}{=}\PH$ compared to $\PH$ (BM25-only) and $\PH$ (\bgem-only),
with a difference that is statistically [significant/not significant,
pending results] under McNemar's test.  The reliability breakdown
(Table~\ref{tab:reliability}) reveals that \PH{} queries constitute
hard cases where first-stage retrieval under the current corpus fails
regardless of retrieval condition, establishing the practical ceiling
for grounding support.  The \PH{} queries where the hybrid strictly
improves over both baselines are consistent with the prediction that
BM25 and \bgem{} fail on disjoint subsets: vocabulary-mismatch queries
that BM25 misses are recovered by semantic matching, while
definitional queries that \bgem{} scatters across many instruments are
anchored by BM25's exact-term precision.

\paragraph{RQ2: Sufficiency for grounded responses.}
The hard-query rate of \PH{} implies that for a fraction of legal
questions, no amount of retrieval fusion will surface the relevant
provision---either because the provision is absent from the 72-instrument
corpus, or because the query resolves a cross-instrument dependency
that a flat retrieval pass cannot resolve.  In a deployed grounded
legal AI system, these queries would correctly trigger an
insufficient-context response from the generator, preserving
answer reliability at the cost of coverage.  The hybrid system
[is/is not, pending results] sufficient as a first stage for
the majority of the 71 query types; a re-ranking stage would be
required to address the hard-query tail.

\paragraph{Comparison with prior work.}
\citet{lrage2024} and \citet{hyparag2024} report similar patterns of
hybrid improvement over single-retriever baselines on legal document
collections, though on different jurisdictions and document types.
The efficiency data (Table~\ref{tab:efficiency}) show that concurrent
execution limits the latency overhead of hybrid retrieval to a factor
of \PH{}over the faster single-retriever baseline, which is acceptable
for interactive legal research tools but warrants profiling in
high-throughput batch settings.

\paragraph{Practical implication.}
At a top-$k_f{=}5$ fusion cutoff, first-stage hybrid retrieval
[does/does not, pending results] provide sufficient context for a
strict-grounding generator (Llama~3.3-70B) to produce citation-supported
answers on the majority of EU Climate Law queries.  Increasing $k_f$
would improve coverage but increase the context window cost for the
language model.  This trade-off should be calibrated against downstream
hallucination metrics in a full end-to-end evaluation, which remains
as future work (\S\ref{sec:conclusion}).

% ──────────────────────────────────────────────────────────────────────────────
\section{Error Analysis}
\label{sec:error}

We examine five representative failure modes.  For each, we state the
query, the expected CELEX instrument(s), hypothesised failure mode,
and one design implication.  Retrieved results are marked \PH{} pending
completion of evaluation runs.

\paragraph{Case~1: Vocabulary mismatch (BM25 failure).}
\textit{Query}: ``What monitoring obligations apply to aircraft operators
under the EU ETS?''
\textit{Expected}: \texttt{32003L0087} (Articles~14--15).
\textit{BM25 top-1}: \PH{} (maritime MRV instrument due to high
term-frequency overlap on \textit{monitoring}, \textit{operators},
\textit{obligations}).
\textit{\bgem{} and hybrid}: \PH.
\textit{Failure mode}: vocabulary mismatch.
\textit{Implication}: domain-specific query expansion for aviation
ETS vocabulary would reduce this class of BM25 failures without
requiring model fine-tuning.

\paragraph{Case~2: Semantic drift (\bgem{} failure).}
\textit{Query}: ``What is the definition of `installation' under the EU ETS?''
\textit{Expected}: \texttt{32003L0087} (Article~3).
\textit{\bgem{} top-5}: \PH{} (definitions articles from CBAM, Energy
Efficiency Directive, and LULUCF, which all define the word
\textit{installation} in their respective scope sections).
\textit{BM25 top-1}: \texttt{32003L0087} (exact-term match on
``installation'' in Article~3 heading).
\textit{Failure mode}: semantic drift.
\textit{Implication}: definitional queries are susceptible to density
across many definitions articles; a structural signal (article-type
metadata) could be used to prioritise definitions-article provisions
in re-ranking.

\paragraph{Case~3: Article fragmentation (all systems fail).}
\textit{Query}: ``What criteria govern the free allocation of allowances
to industrial installations at risk of carbon leakage in ETS Phase~4?''
\textit{Expected}: \texttt{32003L0087} and the Phase~4 benchmarks
regulation.
\textit{All systems}: \PH{} (individual articles from one instrument
appear in top-5, but the cross-instrument answer spans provisions
not simultaneously present).
\textit{Failure mode}: article fragmentation.
\textit{Implication}: multi-hop retrieval across instruments, or an
increased $k_f$, is required to support cross-instrument legal
questions.

\paragraph{Case~4: Cross-reference gap (hybrid partial failure).}
\textit{Query}: ``How does the Carbon Border Adjustment Mechanism
interface with EU ETS Phase~4 for electricity imports?''
\textit{Expected}: \texttt{32023R0956} (CBAM Regulation) and
\texttt{32003L0087} (ETS Directive).
\textit{Hybrid top-5}: \PH{} (CBAM articles dominate; ETS
cross-references at rank~6+, outside $k_f{=}5$).
\textit{Failure mode}: cross-reference gap.
\textit{Implication}: increasing $k_f$ or introducing an explicit
cross-reference expansion step would recover the second relevant
instrument.

\paragraph{Case~5: Definitional query across governance instruments (sparse failure).}
\textit{Query}: ``What reporting period applies to national greenhouse gas
inventories under the European Climate Law?''
\textit{Expected}: \texttt{32021R1119} (European Climate Law).
\textit{BM25 top-1}: \PH{} (MRV Regulation scores highest owing to
repeated co-occurrence of \textit{reporting period} and
\textit{greenhouse gas}).
\textit{\bgem{} and hybrid}: \PH.
\textit{Failure mode}: definitional query with high-frequency competing terms.
\textit{Implication}: governance-level queries (framing, reporting
frameworks) favour dense retrieval; BM25 requires instrument-scoped
search or query rewriting to resolve ambiguity with sector-specific
regulations.

% ──────────────────────────────────────────────────────────────────────────────
\section{Threats to Validity}
\label{sec:validity}

\subsection{Internal Validity}
\label{sec:validity-internal}

The gold benchmark was constructed by the paper's authors, creating
a risk of selection bias toward queries whose structure aligns with
the system's retrieval strengths.  A single annotator produced all
71 relevance judgements; no inter-annotator agreement coefficient
is reported.  We acknowledge this as a limitation: independent
annotation of a 20\,\% sample is required before the benchmark
can be considered suitable for primary comparison in published IR
evaluation campaigns.

\subsection{External Validity}
\label{sec:validity-external}

Results are obtained on a single domain (EU Climate Law) and may not
transfer to other bodies of EU legislation (competition law, GDPR,
financial regulation) or to other jurisdictions.  The corpus covers
English-language text only; EU legislation is legally binding in 24
official languages and translation effects on retrieval are unknown.
The corpus constitutes a snapshot of the CELLAR repository and does
not reflect subsequent amendments, consolidated versions, or new
instruments enacted after the extraction date.  72 CELEX identifiers
cover the core EU climate \textit{acquis} but exclude numerous
delegated and implementing acts.

\subsection{Construct Validity}
\label{sec:validity-construct}

Binary relevance cannot distinguish between highly material and
marginally relevant instruments; two instruments judged relevant for
the same query may differ substantially in how directly they answer it.
CELEX-level matching does not assess whether the retrieved
\emph{article} within the instrument is the most relevant; a retriever
could achieve high Hit Rate by returning the correct instrument at rank
1 via an irrelevant article that happens to share vocabulary.
Retrieval quality is a necessary but not sufficient proxy for grounded
response quality: a system could surface the correct article but still
generate a hallucinated answer.

Regarding context-window truncation: the maximum article in the corpus
(137,713 characters, ${\approx}34\,000$ tokens) far exceeds \bgem's
8,192-token limit.  Sentence-transformers silently truncates these
inputs.  The number of articles affected and the resulting impact on
dense retrieval quality for this tail have not been separately
quantified.

\subsection{Reproducibility}
\label{sec:validity-reproducibility}

All library versions are pinned in \texttt{requirements.txt}.  The
\bgem model is loaded by name (\texttt{BAAI/bge-m3}); the weight hash
should be pinned at download time to guard against upstream model
updates.  All hyperparameters are centralised in \texttt{config.py}
and overridable via environment variables.  The generator (Llama~3.3-70B
via Ollama) is used only for qualitative examples; all quantitative
metrics are retrieval-only and do not depend on the generator.  Code
and benchmark will be released at [repository URL].

% ──────────────────────────────────────────────────────────────────────────────
\section{Conclusion}
\label{sec:conclusion}

\textbf{RQ1.}  \PH{} (pending results; to be answered from
Table~\ref{tab:effectiveness} and Table~\ref{tab:reliability}).

\textbf{RQ2.}  \PH{} (pending results; to be answered from the
hard-query rate in Table~\ref{tab:reliability}).

The strongest result is that [PLACEHOLDER --- to be stated once
experimental runs are complete].  The most important limitation is
that the benchmark is constructed by a single annotator and has not
been independently validated, constraining the generalisability of
reported metrics.

\paragraph{Future work.}
\begin{enumerate}
  \item \textbf{Graded relevance annotation.}  Adding three-level
    relevance judgements (0/1/2) to a subset of the benchmark would
    permit nDCG computation and enable more granular gain analysis.
  \item \textbf{Cross-encoder re-ranking.}  Deploying a cross-encoder
    on the top-5 hybrid output would directly test whether re-ranking
    closes the gap on hard queries.
  \item \textbf{Multilingual extension.}  Extracting parallel French
    and German versions of the same 72 instruments from CELLAR would
    enable the first multilingual EU Climate Law retrieval evaluation,
    exploiting \bgem's multilingual capability.
\end{enumerate}

% ──────────────────────────────────────────────────────────────────────────────
\section{Generative AI Disclosure}
\label{sec:ai-disclosure}

Claude\,Sonnet (Anthropic) was used to assist with code scaffolding
for the retrieval pipeline and for drafting portions of this manuscript.
GitHub Copilot assisted with boilerplate Python.  Llama~3.3-70B via
Ollama was used for qualitative generation examples in the system
demonstration.  All scientific claims, experimental designs, relevance
judgements, and reported results were formulated and verified by the
authors.  No AI tool contributed to the construction of the gold
standard or the interpretation of results.

% ──────────────────────────────────────────────────────────────────────────────

\bibliographystyle{ios1}
\bibliography{shafat2025_references}

\end{document}
